%=============================================================================
% Chapter 15: Deployment
%=============================================================================
\chapter{Deployment}

\section{Deployment Platform: Vercel}

FloodEye is optimised for deployment on Vercel, the cloud platform created by the team behind Next.js. Vercel provides zero-configuration deployment for Next.js applications with automatic serverless function provisioning, global CDN distribution, and Git-integrated continuous deployment.

\section{Deployment Architecture}

\begin{figure}[H]
\centering
\resizebox{\linewidth}{!}{%
\begin{tikzpicture}[node distance=0.8cm]
  \node[startend, text width=3cm] (github) {GitHub\\Repository};
  \node[process, text width=3cm, right=2cm of github] (ci) {Vercel Build\\Pipeline};
  \node[process, text width=3cm, right=2cm of ci] (build) {Next.js Build\\\texttt{--webpack}};
  \node[process, text width=3cm, right=2cm of build] (cdn) {Vercel CDN\\(Global Edge)};

  \node[process, text width=3cm, below=1.5cm of ci] (env) {Environment\\Variables Injected};
  \node[process, text width=3cm, below=1.5cm of build] (pwa) {PWA Assets\\Generated};
  \node[startend, text width=3cm, below=1.5cm of cdn] (user) {End Users\\(Browser / PWA)};

  \draw[arrow] (github) -- node[above,font=\tiny]{push} (ci);
  \draw[arrow] (ci) -- (build);
  \draw[arrow] (build) -- (cdn);
  \draw[arrow] (ci) -- (env);
  \draw[arrowg] (env) -- (build);
  \draw[arrowg] (build) -- (pwa);
  \draw[arrow] (cdn) -- (user);
\end{tikzpicture}%
}
\caption{FloodEye Vercel Deployment Pipeline}
\label{fig:deployment}
\end{figure}

\section{Build Process}

The production build is triggered by:
\begin{lstlisting}[language={},caption={Build command}]
npm run build   -- internally runs: next build --webpack
\end{lstlisting}

The \texttt{--webpack} flag enables the webpack bundler (rather than Turbopack) for the production build, ensuring compatibility with all dependencies. Key build outputs:

\begin{itemize}[leftmargin=1.5em]
  \item \textbf{Server functions:} Next.js Route Handlers compiled to Vercel serverless functions.
  \item \textbf{Static assets:} JavaScript bundles, CSS, and public files distributed to Vercel's global CDN.
  \item \textbf{PWA service worker:} \texttt{next-pwa} (Workbox) generates \texttt{sw.\allowbreak{}js} and \texttt{workbox-*.\allowbreak{}js} during build. These are excluded in development (\texttt{disable: process.\allowbreak{}env.\allowbreak{}NODE\_\allowbreak{}ENV === 'development'}).
\end{itemize}

\section{Environment Configuration}

Environment variables are configured in the Vercel project dashboard (not in the repository):

\begin{table}[H]
\centering
\caption{Vercel Environment Variable Configuration}
\begin{tabularx}{\linewidth}{|l|l|X|}
\hline
\rowcolor{FloodLight}
\textbf{Variable} & \textbf{Scope} & \textbf{Notes} \\
\hline
\texttt{THINGSPEAK\_\allowbreak{}CHANNEL\_\allowbreak{}ID} & Server & Never exposed to the browser \\
\texttt{THINGSPEAK\_\allowbreak{}READ\_\allowbreak{}API\_\allowbreak{}KEY} & Server & Read-only; no write risk \\
\texttt{NEXT\_\allowbreak{}PUBLIC\_\allowbreak{}MAPTILER\_\allowbreak{}KEY} & Client & Prefixed \texttt{NEXT\_\allowbreak{}PUBLIC\_\allowbreak{}}; included in browser bundle \\
\hline
\end{tabularx}
\end{table}

\section{PWA Configuration}

PWA is configured in \texttt{next.\allowbreak{}config.\allowbreak{}ts} using \texttt{@ducanh2912/\allowbreak{}next-pwa}:

\begin{lstlisting}[language=TypeScript,caption={PWA configuration in next.config.ts}]
const withPWA = withPWAInit({
  dest: "public",
  cacheOnFrontEndNav: true,
  aggressiveFrontEndNavCaching: true,
  reloadOnOnline: true,
  disable: process.env.NODE_ENV === "development",
  workboxOptions: { disableDevLogs: true },
});
\end{lstlisting}

The PWA manifest is defined in \texttt{app/\allowbreak{}manifest.\allowbreak{}ts} and specifies the app name, icons (\texttt{icon-192x192.\allowbreak{}png}, \texttt{icon-512x512.\allowbreak{}png}), theme colour, background colour, and display mode (\texttt{standalone}).

\section{Deployment Steps}

\begin{enumerate}[leftmargin=1.5em]
  \item Push code to the GitHub repository (main branch).
  \item Vercel detects the push and triggers a new build.
  \item Environment variables are injected from the Vercel dashboard.
  \item Next.js build runs (\texttt{next build --webpack}).
  \item Build outputs are deployed to Vercel's global edge network.
  \item The new deployment is promoted to the production URL (e.g.\ \texttt{floodeye.\allowbreak{}vercel.\allowbreak{}app}).
  \item Preview deployments are automatically created for every pull request.
\end{enumerate}

\section{Local Development}

\begin{lstlisting}[language={},caption={Local development startup}]
# 1. Clone and install
git clone <repo-url>
cd IOT-TIH
npm install

# 2. Create .env.local
echo "THINGSPEAK_CHANNEL_ID=your_id" >> .env.local
echo "THINGSPEAK_READ_API_KEY=your_key" >> .env.local
echo "NEXT_PUBLIC_MAPTILER_KEY=your_key" >> .env.local

# 3. Start dev server
npm run dev
# Opens http://localhost:3000
\end{lstlisting}

\section{Performance Characteristics}

\begin{itemize}[leftmargin=1.5em]
  \item \textbf{Cold start:} Vercel serverless functions have a cold start of $\sim$100--300 ms for the first request after a period of inactivity.
  \item \textbf{Cached API response:} $<$ 5 ms (served from CDN edge or server memory).
  \item \textbf{Live ThingSpeak fetch:} $\sim$100--200 ms (HTTP round-trip to ThingSpeak API).
  \item \textbf{ML model load:} $\sim$1--3 s (fetching and deserialising 291 KB binary weights + model JSON).
  \item \textbf{ML inference:} $<$ 50 ms per prediction (WebGL-accelerated on modern devices).
\end{itemize}
